\chapter{Combinatorial}

\section{Permutations}
	\subsection{Factorial}
		% \import{factorial.tex}
		\kactlimport{IntPerm.h}

	\subsection{Cycles}
		Let $g_S(n)$ be the number of $n$-permutations whose cycle lengths all belong to the set $S$. Then
		$$\sum_{n=0} ^\infty g_S(n) \frac{x^n}{n!} = \exp\left(\sum_{n\in S} \frac{x^n} {n} \right)$$

	\subsection{Derangements}
		Permutations of a set such that none of the elements appear in their original position.
		\[ \mkern-2mu D(n) = (n-1)(D(n-1)+D(n-2)) = n D(n-1)+(-1)^n = \left\lfloor\frac{n!}{e}\right\rceil \]

	\subsection{Burnside's lemma}
		Given a group $G$ of symmetries and a set $X$, the number of elements of $X$ \emph{up to symmetry} equals
		 \[ {\frac {1}{|G|}}\sum _{{g\in G}}|X^{g}|, \]
		 where $X^{g}$ are the elements fixed by $g$ ($g.x = x$).

		 If $f(n)$ counts ``configurations'' (of some sort) of length $n$, we can ignore rotational symmetry using $G = \mathbb Z_n$ to get
		 \[ g(n) = \frac 1 n \sum_{k=0}^{n-1}{f(\text{gcd}(n, k))} = \frac 1 n \sum_{k|n}{f(k)\phi(n/k)}. \]

\section{Partitions and subsets}
	\subsection{Partition function}
		Number of ways of writing $n$ as a sum of positive integers, disregarding the order of the summands.
		\[ p(0) = 1,\ p(n) = \sum_{k \in \mathbb Z \setminus \{0\}}{(-1)^{k+1} p(n - k(3k-1) / 2)} \]
		% \[ p(n) \sim 0.145 / n \cdot \exp(2.56 \sqrt{n}) \]

		% \begin{center}
		% \begin{tabular}{c|c@{\ }c@{\ }c@{\ }c@{\ }c@{\ }c@{\ }c@{\ }c@{\ }c@{\ }c@{\ }c@{\ }c@{\ }c}
		% 	$n$    & 0 & 1 & 2 & 3 & 4 & 5 & 6  & 7  & 8  & 9  & 20  & 50  & 100 \\ \hline
		% 	$p(n)$ & 1 & 1 & 2 & 3 & 5 & 7 & 11 & 15 & 22 & 30 & 627 & $\mathtt{\sim}$2e5 & $\mathtt{\sim}$2e8 \\
		% \end{tabular}
		% \end{center}

	% \subsection{Binomials}
	% 	\kactlimport{multinomial.h}

	\subsection{Lucas' Theorem}
		Let $n,m$ be non-negative integers and $p$ a prime. Write $n=n_kp^k+...+n_1p+n_0$ and $m=m_kp^k+...+m_1p+m_0$. Then $\binom{n}{m} \equiv \prod_{i=0}^k\binom{n_i}{m_i} \pmod{p}$.

	\subsection{Hockey-stick Identity}
		\[ \sum_{n=a}^b \binom{n}{m} = \binom{b+1}{m+1} - \binom{a}{m+1} \]
		
		Let $f_z=\sum_{i=0}^x \binom{ki+y}{z}$ for fixed $x, y, k$, can compute $f_i$ for all $i \leq m$ in $O(km)$ time using recurrence ($f_0=x+1$):

		$$ f_z = \frac{1}{k}(\binom{k(x+1)+y}{z+1} - \binom{y}{z+1} - \sum_{j=2}^{\min(k,z+1)} \binom{k}{j} f_{z-j+1}) $$

	\subsection{Vandermonde's Identity}
		\[ \sum_{k=0}^r \binom{m}{k}\binom{n}{r-k} = \binom{n+m}{r} \]
		
\section{Generating Functions}
	OGF of $a_n$ is $A(x)=\sum_{n \ge 0}a_nx^n$. EGF is $\sum_{n \ge 0}\frac{a_n}{n!}x^n$.

	\subsection{Operations and Common Tricks}

	In recurrence equations, multiply a power of $x$ to both sides and summing up for all $n$. 

	Multiply by $n$ ($c_n=na_n$), OGF/EGF: $C(x)=xA'(x)$.

	Shifting $c_n = a_{n-k}$ where $a_{<0}=0$, OGF: $C(x)=x^kA(x)$; EGF: Integrate $k$ times.

	Shifting $c_n = a_{n+k}$, OGF: $C(x)=\frac{A(x)-(a_0+a_1x+\ldots+a_{k-1}x^{k-1})}{x^k}$; EGF: Differentiate $k$ times.

	Taylor shift $C(x)=A(x+k)$ assuming $A(x)$ has degree $n$, OGF: $c_j=\frac{[x^{n-j}](B(x)D(x))}{j!}$, $[x^i]B(x)=a_{n-i}(n-i)!$, $[x^i]D(x)=\frac{k^i}{i!}$, $B(x), C(x), D(x)$ all have degree $n$; EGF: $c_j=[x^{n-j}](B(x)D(x))$, $[x^i]B(x)=a_{n-i}$, $[x^i]D(x)=\frac{k^i}{i!}$.

	Power of $k$, $C(x) = A(x)^k$, OGF: $c_n=\sum_{i_1+\ldots+i_k=n}a_{i_1}\cdots a_{i_k}$; 
	EGF: $c_n=\sum_{i_1+\ldots+i_k=n}\frac{n!}{i_1!\cdots i_k!}a_{i_1}\cdots a_{i_k}$

	If $a_n=$ number of ways to make a single component using $n$ labelled items and has EGF $A(x)$, then number of ways to make several unordered components using $n$ items has EGF $e^{A(x)}$. (Fix number of components to be $k$, then it's $\frac{A(x)^k}{k!}$)

	Prefix Sum Trick, OGF only: $c_n = a_0+\ldots+a_n$, $C(x)=\frac{1}{1-x}A(x)$.

	Newton's Method: to solve $F(Q(x))=0$ for some poly-valued function $F$, $Q_{k+1}=Q_k-\frac{F(Q_k)}{F'(Q_k)} \mod x^{2^{k+1}}$.

	Reciprocal of $A(x)$ is $Q(x)$: (if division by zero is an issue, multiply both sides of the fraction by something first) let $F(Q)=\frac{1}{Q}-A$, then $Q_{k+1}=Q_k(2-AQ_k) \mod x^{2^{k+1}}$. Require $[x^0]A(x) \neq 0$.

	Logarithm: $(\ln A(x))'=\frac{A'(x)}{A(x)}$, integrate both sides. Usually require $[x^0]A(x)=1$.

	Exponent: let $F(Q)=\ln Q-A$, then $Q_{k+1}=Q_k(1+A-\ln Q_k) \mod x^{2^{k+1}}$. Usually require $[x^0]A(x)=0$.

	\subsection{Useful Equations and Applications}

	$$ \frac{1}{1-x} = 1+x+x^2+\ldots = \sum_{n \ge 0} x^n \ ;\ \frac{1}{(1-x)^k} = \sum_{n \ge 0} \binom{n+k-1}{n} x^n $$

	For different $x_i$: (expand $H(t)=\frac{1}{(1-x_1 t)\cdots (1-x_n t)}$)

	$$ \sum_{a_1+\ldots+a_n=m} \Pi_{i=1}^n x_i ^ {a_i} = \sum_{i=1}^n \frac{x_i ^ {n+m-1}}{\Pi_{j \neq i} (x_i-x_j)} $$

	$\sum_{i \bmod m = r}\binom{n}{i}=[x^0]F(x)$, where $F(x)=x^{m-r}(1+x)^n$ modulo $x^m-1$ ($x^m=1$, $x^{m+1}=x$, \ldots) can be found in $O(m^2 \log n)$ using binary exponentiation and multiplying polys naively.

	Partition of $n$ to odd parts = Partition to distinct parts

	$P[$stopping time $= k] = \sum P[$got to stopping state in $i$ moves$]\cdot P[$stay at stopping state in $k-i$ moves$]$

\section{General purpose numbers}
	\subsection{Bernoulli numbers}
		EGF of Bernoulli numbers is $B(t)=\frac{t}{e^t-1}$ (FFT-able).
		$B[0,\ldots] = [1, -\frac{1}{2}, \frac{1}{6}, 0, -\frac{1}{30}, 0, \frac{1}{42}, \ldots]$

		Sums of powers:
		\small
		\[ \sum_{i=1}^n n^m = \frac{1}{m+1} \sum_{k=0}^m \binom{m+1}{k} B_k \cdot (n+1)^{m+1-k} \]
		\normalsize

		Euler-Maclaurin formula for infinite sums:
		\small
		\[ \sum_{i=m}^{\infty} f(i) = \int_m^\infty f(x) dx - \sum_{k=1}^\infty \frac{B_k}{k!}f^{(k-1)}(m) \]
		\[ \approx \int_{m}^\infty f(x)dx + \frac{f(m)}{2} - \frac{f'(m)}{12} + \frac{f'''(m)}{720} + O(f^{(5)}(m)) \]
		\normalsize

	\subsection{Stirling numbers of the first kind}
		Number of permutations on $n$ items with $k$ cycles is $|s(n,k)|$.
		$s(n,k)=(-1)^{n-k}|s(n,k)|$.
		\begin{align*}
			&|s(n,k)| = |s(n-1,k-1)| + (n-1) |s(n-1,k)|,\ s(0,0) = 1 \\
			&\textstyle (x)_n = x(x-1) \dots (x-n+1) = \sum_{k=0}^n s(n,k)x^k \\
			&\textstyle x^{(n)} = x(x+1) \dots (x+n-1) = (-1)^n(-x)_n = \sum_{k=0}^n |s(n,k)|x^k
		\end{align*}
		For fixed $k$, $s(n,k) = [\frac{x^n}{n!}]\frac{\log^k (1+x)}{k!}$. 
		For fixed $n$, divide and conquer and expand $(x)_n$ in $O(n \log^2 n)$. 
		For single log, note $(x)_{2n}=(x)_n(x-n)_n$, so first find $(x)_n$, 
		use Taylor shift for $(x-n)_n$, then multiply them.
		$|s(8,k)| = 8, 0, 5040, 13068, 13132, 6769, 1960, 322, 28, 1$ \\
		$|s(n,2)| = 0, 0, 1, 3, 11, 50, 274, 1764, 13068, 109584, \dots$

	\subsection{Eulerian numbers}
		Number of permutations $\pi \in S_n$ in which exactly $k$ elements are greater than the previous element. $k$ $j$:s s.t. $\pi(j)>\pi(j+1)$, $k+1$ $j$:s s.t. $\pi(j)\geq j$, $k$ $j$:s s.t. $\pi(j)>j$.
		$$E(n,k) = (n-k)E(n-1,k-1) + (k+1)E(n-1,k)$$
		$$E(n,0) = E(n,n-1) = 1$$
		$$E(n,k) = \sum_{j=0}^k(-1)^j\binom{n+1}{j}(k+1-j)^n$$

	\subsection{Stirling numbers of the second kind}
		Partitions of $n$ distinct elements into exactly $k$ groups.
		$$S(n,k) = S(n-1,k-1) + k S(n-1,k),\ S(n,1) = S(n,n) = 1$$
		$$S(n,k) = \frac{1}{k!}\sum_{j=0}^k (-1)^{k-j}\binom{k}{j}j^n,\ x^n = \sum_{k=0}^n S(n,k)(x)_k $$
		For fixed $k$, $S(n,k) = [\frac{x^n}{n!}]\frac{(e^x-1)^k}{k!}$.
		For fixed $n$, $S(n,k)=\sum_{j=0}^k \frac{(-1)^{k-j}}{(k-j)!} \cdot \frac{j^n}{j!}$ (compute by FFT).
	
	\subsection{Subset Inversion}
		$f(S)$ is the answer for set $S$. $g$ is easy to compute $\implies$ recover $f$ with PIE. 
		Can also change summmation to $T:T\subseteq S$.

		$$g(S)=\sum_{T:S\subseteq T}f(T) \iff f(S)=\sum_{T:S\subseteq T}(-1)^{|T|-|S|}g(T)$$

		Stirling Inversion applies when $g(n)$ is counting $n$ things with overlap allowed, where overlap is not allowed in $f(n)$. 
		We can compute $g$ from $f$ by distribute $n$ items into $k$ buckets then count with $k$ distinct elements, 
		but also in reverse.

		$$g(n)=\sum_{k=0}^nS(n,k)f(k) \iff f(n)=\sum_{k=0}^n s(n,k)g(k)$$

		Mobius min/max Inversion (min+max can be swapped or replaced with gcd+lcm or AND+OR).

		$$\min(S) = \sum_{\emptyset \neq T \subseteq S}(-1)^{|T|+1}\max(T)$$

	\subsection{Bell numbers}
		Total number of partitions of $n$ distinct elements. $B(n) =$
		$1, 1, 2, 5, 15, 52, 203, 877, 4140, 21147, \dots$. For $p$ prime,
		\[ B(p^m+n)\equiv mB(n)+B(n+1) \pmod{p} \]

	\subsection{Labeled unrooted trees}
		\# on $n$ vertices: $n^{n-2}$ \\
		\# on $k$ existing trees of size $n_i$: $n_1n_2\cdots n_k n^{k-2}$ \\
		\# with degrees $d_i$: $(n-2)! / ((d_1-1)! \cdots (d_n-1)!)$

	\subsection{Catalan numbers}
		\[ C_n=\frac{1}{n+1}\binom{2n}{n}= \binom{2n}{n}-\binom{2n}{n+1} = \frac{(2n)!}{(n+1)!n!} \]
		\[ C_0=1,\ C_{n+1} = \frac{2(2n+1)}{n+2}C_n,\ C_{n+1}=\sum C_iC_{n-i} \]
		${C_n = 1, 1, 2, 5, 14, 42, 132, 429, 1430, 4862, 16796, 58786, \dots}$
		\begin{itemize}[noitemsep]
			\item sub-diagonal monotone paths in an $n\times n$ grid.
			\item strings with $n$ pairs of parenthesis, correctly nested.
			\item binary trees with with $n+1$ leaves (0 or 2 children).
			\item ordered trees with $n+1$ vertices.
			\item ways a convex polygon with $n+2$ sides can be cut into triangles by connecting vertices with straight lines.
			\item permutations of $[n]$ with no 3-term increasing subseq.
		\end{itemize}
